\section{Límites proyectivos en álgebra homológica}
\label{section:0.13}

\subsection{La condición de Mittag-Leffler}
\label{subsection:0.13.1}

\begin{env}[13.1.1]
\label{0.13.1.1}
Sea $\cat{C}$ una categoría abeliana en la que existan los productos
infinitos (axioma AB~3* de T, 1.5).
Entonces existe el ínfimo de una familia de subobjetos de un objeto de
$\cat{C}$, y todo sistema proyectivo de objetos de $\cat{C}$ posee un
límite proyectivo; este último es un funtor exacto por la izquierda del
sistema proyectivo en cuestión (T, 1.8).
\oldpage[0\textsubscript{III}]{65}
Sea $(A_\alpha,f_{\alpha\beta})$ un sistema proyectivo de objetos de
$\cat{C}$ cuyo conjunto de índices $\mathrm{I}$ está filtrado por la derecha.
Pongamos
\[
  A=\varprojlim A_\alpha,
\]
y, para cada $\alpha\in\mathrm{I}$, sea
$f_\alpha:A\longrightarrow A_\alpha$ el morfismo canónico.
Para todo $\alpha\in\mathrm{I}$, los subobjetos
$f_{\alpha\beta}(A_\beta)$, para $\alpha\leq\beta$, forman una familia
filtrada decreciente de subobjetos de $A_\alpha$.
El subobjeto
\[
  A'_\alpha
  =
  \inf_{\beta\geq\alpha}f_{\alpha\beta}(A_\beta)
\]
se llama el subobjeto de las \emph{imágenes universales} en $A_\alpha$.
Está claro que
\[
  f_\alpha(A)\subset A'_\alpha,
  \qquad
  f_{\alpha\beta}(A'_\beta)\subset A'_\alpha
  \quad(\alpha\leq\beta).
\]
Así,
$(A'_\alpha,f_{\alpha\beta}|A'_\beta)$ es un sistema proyectivo y
\[
  A=\varprojlim A'_\alpha.
\]
\end{env}

\begin{env}[13.1.2]
\label{0.13.1.2}
Para un sistema proyectivo $(A_\alpha,f_{\alpha\beta})$ en $\cat{C}$,
se llama \emph{condición de Mittag-Leffler} a la siguiente:
\begin{itemize}
  \item[\textbf{(ML)}]
  \emph{Para todo índice $\alpha$, existe un $\beta\geq\alpha$ tal
  que, para todo $\gamma\geq\beta$,}
  \[
    f_{\alpha\gamma}(A_\gamma)
    =
    f_{\alpha\beta}(A_\beta).
  \]
\end{itemize}

Está claro que la condición (ML) se cumple si todos los
$f_{\alpha\beta}$ son epimorfismos.
Recíprocamente, supongamos que se cumple (ML), y para cada
$\alpha\in\mathrm{I}$ sea $A'_\alpha$ el subobjeto de las imágenes
universales en $A_\alpha$.
Entonces, para $\alpha\leq\beta$, la restricción de
$f_{\alpha\beta}$ a $A'_\beta$ es un \emph{epimorfismo}
\[
  A'_\beta\longrightarrow A'_\alpha.
\]
En efecto, elijamos $\gamma\geq\beta$ tal que
\[
  f_{\beta\delta}(A_\delta)
  =
  f_{\beta\gamma}(A_\gamma)
  \qquad(\delta\geq\gamma).
\]
Entonces $A'_\beta=f_{\beta\gamma}(A_\gamma)$, y por tanto
\[
  f_{\alpha\beta}(A'_\beta)
  =
  f_{\alpha\gamma}(A_\gamma),
  \qquad
  A'_\alpha
  =
  f_{\alpha\gamma}(A_\gamma)
  =
  f_{\alpha\beta}(A'_\beta).
\]

La condición (ML) se cumple también cuando los objetos $A_\alpha$ son
\emph{artinianos} en $\cat{C}$, es decir, cuando toda familia de
subobjetos de $A_\alpha$ posee un elemento minimal.
En efecto, un elemento minimal de la familia filtrada decreciente
$\bigl(f_{\alpha\beta}(A_\beta)\bigr)$ de subobjetos de $A_\alpha$
es necesariamente el menor de esos subobjetos.
\end{env}

\begin{env}[13.1.3]
\label{0.13.1.3}
\emph{Observación} (13.1.3).
La condición (ML) puede formularse también cuando $\cat{C}$ es, por
ejemplo, la categoría de conjuntos.
Se puede definir de nuevo el subconjunto de las imágenes universales en
$A_\alpha$, y las observaciones hechas en \sref{0.13.1.1} y
\sref{0.13.1.2} siguen siendo válidas.
\end{env}

\subsection{La condición de Mittag-Leffler para grupos abelianos}
\label{subsection:0.13.2}

\begin{proposition}[13.2.1]
\label{0.13.2.1}
Sea
\[
  0\longrightarrow A_\alpha
  \xrightarrow{u_\alpha}B_\alpha
  \xrightarrow{v_\alpha}C_\alpha
  \longrightarrow0
\]
una sucesión exacta de sistemas proyectivos de grupos abelianos, relativos
al mismo conjunto de índices filtrado $\mathrm{I}$.
\begin{enumerate}
  \item[{\rm(i)}]
  Si $(B_\alpha)$ satisface {\rm(ML)}, entonces también $(C_\alpha)$.
  \item[{\rm(ii)}]
  Si $(A_\alpha)$ y $(C_\alpha)$ satisfacen {\rm(ML)}, entonces también
  $(B_\alpha)$.
\end{enumerate}
\end{proposition}

\begin{proof}
Sean $(f_{\alpha\beta})$, $(g_{\alpha\beta})$ y
$(h_{\alpha\beta})$ los sistemas de homomorfismos que definen
$(A_\alpha)$, $(B_\alpha)$ y $(C_\alpha)$, respectivamente.

(i) Supongamos que
\[
  g_{\alpha\beta}(B_\beta)
  =
  g_{\alpha\lambda}(B_\lambda)
  \qquad(\lambda\geq\beta).
\]
Como $v_\beta$ y $v_\lambda$ son sobreyectivos,
\[
  h_{\alpha\beta}(C_\beta)
  =
  v_\alpha\bigl(g_{\alpha\beta}(B_\beta)\bigr)
  =
  v_\alpha\bigl(g_{\alpha\lambda}(B_\lambda)\bigr)
  =
  h_{\alpha\lambda}(C_\lambda)
\]
para $\lambda\geq\beta$.

(ii) Sea $\alpha\in\mathrm{I}$.
Elijamos $\beta\geq\alpha$ tal que
\[
  f_{\alpha\beta}(A_\beta)
  =
  f_{\alpha\lambda}(A_\lambda)
  \qquad(\lambda\geq\beta),
\]
y elijamos después $\gamma\geq\beta$ tal que
\[
  h_{\beta\gamma}(C_\gamma)
  =
  h_{\beta\lambda}(C_\lambda)
  \qquad(\lambda\geq\gamma).
\]
Sea $y_\alpha\in g_{\alpha\gamma}(B_\gamma)$, escribamos
$y_\alpha=g_{\alpha\gamma}(y_\gamma)$ con $y_\gamma\in B_\gamma$,
y pongamos $y_\beta=g_{\beta\gamma}(y_\gamma)$.
Entonces
\[
  v_\beta(y_\beta)
  =
  h_{\beta\gamma}\bigl(v_\gamma(y_\gamma)\bigr).
\]
Para todo $\lambda\geq\gamma$, existe, por hipótesis, un
$y_\lambda\in B_\lambda$ tal que
\[
  h_{\beta\gamma}\bigl(v_\gamma(y_\gamma)\bigr)
  =
  h_{\beta\lambda}\bigl(v_\lambda(y_\lambda)\bigr).
\]
De aquí que
$v_\beta(y_\beta-g_{\beta\lambda}(y_\lambda))=0$, y por tanto
existe un $x_\beta\in A_\beta$
\oldpage[0\textsubscript{III}]{66}
tal que
\[
  y_\beta
  =
  g_{\beta\lambda}(y_\lambda)+u_\beta(x_\beta).
\]
Se sigue que
\[
  y_\alpha
  =
  g_{\alpha\lambda}(y_\lambda)
  +
  u_\alpha\bigl(f_{\alpha\beta}(x_\beta)\bigr).
\]
Como $\lambda\geq\beta$, existe un $x_\lambda\in A_\lambda$ tal
que
\[
  f_{\alpha\beta}(x_\beta)=f_{\alpha\lambda}(x_\lambda).
\]
Por consiguiente,
\[
  y_\alpha
  =
  g_{\alpha\lambda}
  \bigl(y_\lambda+u_\lambda(x_\lambda)\bigr)
  \in g_{\alpha\lambda}(B_\lambda),
\]
lo que demuestra la aserción.
\end{proof}

\begin{proposition}[13.2.2]
\label{0.13.2.2}
Sea $\mathrm{I}$ un conjunto ordenado filtrado que posee un subconjunto
cofinal numerable, y sea
\[
  0\longrightarrow A_\alpha
  \xrightarrow{u_\alpha}B_\alpha
  \xrightarrow{v_\alpha}C_\alpha
  \longrightarrow0
\]
una sucesión exacta de sistemas proyectivos de grupos abelianos indexados
por $\mathrm{I}$.
Si $(A_\alpha)$ satisface {\rm(ML)}, entonces la sucesión
\[
  0\longrightarrow\varprojlim A_\alpha
  \longrightarrow\varprojlim B_\alpha
  \longrightarrow\varprojlim C_\alpha
  \longrightarrow0
\]
es exacta.
\end{proposition}

\begin{proof}
Solo queda demostrar que el homomorfismo
\[
  v=\varprojlim v_\alpha:
  \varprojlim B_\alpha\longrightarrow\varprojlim C_\alpha
\]
es sobreyectivo.
Sea $z=(z_\alpha)\in\varprojlim C_\alpha$, y pongamos
\[
  E_\alpha=v_\alpha^{-1}(z_\alpha).
\]
Los $E_\alpha$ forman manifiestamente un sistema proyectivo de conjuntos no
vacíos mediante las restricciones de los homomorfismos
$g_{\alpha\beta}:B_\beta\longrightarrow B_\alpha$.

Demostremos que este sistema proyectivo satisface (ML).
Identificando $A_\alpha$ con un subgrupo de $B_\alpha$ por medio de
$u_\alpha$, elijamos, para todo $\alpha\in\mathrm{I}$, un índice
$\beta\geq\alpha$ tal que
\[
  g_{\alpha\beta}(A_\beta)
  =
  g_{\alpha\lambda}(A_\lambda)
  \qquad(\lambda\geq\beta).
\]
Afirmamos que también
\[
  g_{\alpha\beta}(E_\beta)
  =
  g_{\alpha\lambda}(E_\lambda)
  \qquad(\lambda\geq\beta).
\]
En efecto, tomemos $y_\lambda\in E_\lambda$ y pongamos
\[
  y_\beta=g_{\beta\lambda}(y_\lambda),
  \qquad
  y_\alpha=g_{\alpha\lambda}(y_\lambda).
\]
Sea $y'_\alpha=g_{\alpha\beta}(y'_\beta)$ para algún
$y'_\beta\in E_\beta$.
Entonces $y'_\beta-y_\beta=x_\beta\in A_\beta$, y por hipótesis existe
un $x_\lambda\in A_\lambda$ tal que
\[
  g_{\alpha\beta}(x_\beta)
  =
  g_{\alpha\lambda}(x_\lambda).
\]
Así,
\[
  y'_\alpha
  =
  g_{\alpha\beta}(y_\beta)+g_{\alpha\beta}(x_\beta)
  =
  g_{\alpha\lambda}(y_\lambda+x_\lambda)
  \in g_{\alpha\lambda}(E_\lambda),
\]
lo que demuestra la afirmación.

Se sabe (Bourbaki, \emph{Top.\ g\'{e}n.}, cap.~II,
\textsection3, teorema~1) que, bajo las hipótesis impuestas a
$\mathrm{I}$, un sistema proyectivo de conjuntos no vacíos que satisface (ML)
posee un límite proyectivo no vacío.
Existe, por tanto, un punto
$y=(y_\alpha)\in\varprojlim E_\alpha$.
Como $v_\alpha(y_\alpha)=z_\alpha$ para todo $\alpha$, se tiene
$z=v(y)$.
\end{proof}

\begin{proposition}[13.2.3]
\label{0.13.2.3}
Bajo las hipótesis sobre $\mathrm{I}$ de \sref{0.13.2.2}, sea
$(\mathrm{K}^\bullet_\alpha)_{\alpha\in\mathrm{I}}$ un sistema proyectivo
de complejos de grupos abelianos
\[
  \mathrm{K}^\bullet_\alpha
  =
  (\mathrm{K}^n_\alpha)_{n\in\bb{Z}},
\]
cuyas diferenciales tienen grado $+1$.
Para todo $n$, existe un homomorfismo canónico funtorial
\begin{equation}
\label{0.13.2.3.1}
\tag{13.2.3.1}
  h_n:
  \HH^n\bigl(\varprojlim\mathrm{K}^\bullet_\alpha\bigr)
  \longrightarrow
  \varprojlim\HH^n(\mathrm{K}^\bullet_\alpha).
\end{equation}
Si, para todo grado $n$, el sistema proyectivo de grupos abelianos
$(\mathrm{K}^n_\alpha)_{\alpha\in\mathrm{I}}$ satisface {\rm(ML)},
entonces todo $h_n$ es sobreyectivo.
Si, además, el sistema proyectivo
$(\HH^{n-1}(\mathrm{K}^\bullet_\alpha))_{\alpha\in\mathrm{I}}$
satisface {\rm(ML)}, entonces $h_n$ es biyectivo.
\end{proposition}

\begin{proof}
Pongamos
\[
  \mathrm{K}^n=\varprojlim\mathrm{K}^n_\alpha
  \qquad(n\in\bb{Z}).
\]
Los homomorfismos $h_n$ están definidos por los diagramas conmutativos
\[
  \xymatrix{
    \cdots\ar[r]
      & \mathrm{K}^{n-1}\ar[r]\ar[d]
      & \mathrm{K}^n\ar[r]\ar[d]
      & \mathrm{K}^{n+1}\ar[r]\ar[d]
      & \cdots\\
    \cdots\ar[r]
      & \mathrm{K}^{n-1}_\alpha\ar[r]
      & \mathrm{K}^n_\alpha\ar[r]
      & \mathrm{K}^{n+1}_\alpha\ar[r]
      & \cdots,
  }
\]
\oldpage[0\textsubscript{III}]{67}
siendo las diferenciales en $\mathrm{K}^\bullet$ los límites proyectivos
de las diferenciales correspondientes en los
$\mathrm{K}^\bullet_\alpha$.
Consideremos las sucesiones exactas
\[
  (*_n)\qquad
  0\longrightarrow
  \mathrm{B}^n(\mathrm{K}^\bullet_\alpha)
  \longrightarrow
  \mathrm{Z}^n(\mathrm{K}^\bullet_\alpha)
  \longrightarrow
  \HH^n(\mathrm{K}^\bullet_\alpha)
  \longrightarrow0
\]
y
\[
  (**_n)\qquad
  0\longrightarrow
  \mathrm{Z}^{n-1}(\mathrm{K}^\bullet_\alpha)
  \longrightarrow
  \mathrm{K}^{n-1}_\alpha
  \longrightarrow
  \mathrm{B}^n(\mathrm{K}^\bullet_\alpha)
  \longrightarrow0.
\]
La hipótesis y la Proposición~\sref{0.13.2.1}~(i) muestran que el
sistema proyectivo
$(\mathrm{B}^n(\mathrm{K}^\bullet_\alpha))_{\alpha\in\mathrm{I}}$
satisface (ML) para todo $n$.
Se sigue de \sref{0.13.2.2} que
\[
  (***_n)\qquad
  0\longrightarrow
  \varprojlim_\alpha\mathrm{B}^n(\mathrm{K}^\bullet_\alpha)
  \longrightarrow
  \varprojlim_\alpha\mathrm{Z}^n(\mathrm{K}^\bullet_\alpha)
  \longrightarrow
  \varprojlim_\alpha\HH^n(\mathrm{K}^\bullet_\alpha)
  \longrightarrow0
\]
es exacta.
Ahora bien,
$\varprojlim_\alpha\mathrm{B}^n(\mathrm{K}^\bullet_\alpha)$
se identifica con un subgrupo de $\mathrm{K}^{n+1}$ que contiene
$\mathrm{B}^n(\mathrm{K}^\bullet)$, mientras que
$\varprojlim_\alpha\mathrm{Z}^n(\mathrm{K}^\bullet_\alpha)$
se identifica con un subgrupo de $\mathrm{Z}^n(\mathrm{K}^\bullet)$.
Por consiguiente, $h_n$ es sobreyectivo.

Supongamos ahora, además, que
$(\HH^{n-1}(\mathrm{K}^\bullet_\alpha))_{\alpha\in\mathrm{I}}$
satisface (ML).
Las sucesiones exactas $(*_{n-1})$ y la
Proposición~\sref{0.13.2.1}~(ii) muestran que
$(\mathrm{Z}^{n-1}(\mathrm{K}^\bullet_\alpha))_{\alpha\in\mathrm{I}}$
satisface (ML).
Aplicando \sref{0.13.2.2} a las sucesiones exactas $(**_n)$ se obtiene una
sucesión exacta
\[
  0\longrightarrow
  \varprojlim_\alpha
    \mathrm{Z}^{n-1}(\mathrm{K}^\bullet_\alpha)
  \longrightarrow
  \mathrm{K}^{n-1}
  \xrightarrow{u}
  \varprojlim_\alpha
    \mathrm{B}^n(\mathrm{K}^\bullet_\alpha)
  \longrightarrow0.
\]
Como
$\varprojlim_\alpha\mathrm{B}^n(\mathrm{K}^\bullet_\alpha)
\supset\mathrm{B}^n(\mathrm{K}^\bullet)$
y el compuesto de $u$ con la inyección
\[
  \varprojlim_\alpha\mathrm{B}^n(\mathrm{K}^\bullet_\alpha)
  \longrightarrow\mathrm{K}^n
\]
es la diferencial
$\mathrm{K}^{n-1}\longrightarrow\mathrm{K}^n$, la sobreyectividad de
$u$ implica
\[
  \varprojlim_\alpha\mathrm{B}^n(\mathrm{K}^\bullet_\alpha)
  =
  \mathrm{B}^n(\mathrm{K}^\bullet).
\]
Así, $h_n$ es inyectivo.
\end{proof}

\begin{env}[13.2.4]
\label{0.13.2.4}
\emph{Observaciones} (13.2.4).
\begin{enumerate}
  \item[{\rm(i)}]
  El argumento de \sref{0.13.2.2}
  (cf.\ Bourbaki, \emph{loc.\ cit.}) muestra que la conclusión de
  esa proposición sigue siendo válida si se supone únicamente que los
  $A_\alpha$ pueden dotarse de estructuras de espacios metrizables completos
  en los que las traslaciones son homeomorfismos, que las aplicaciones
  $f_{\alpha\beta}:A_\beta\longrightarrow A_\alpha$ que definen el
  sistema proyectivo son uniformemente continuas para las métricas en
  cuestión, y que $(A_\alpha)$ satisface la condición
  \begin{itemize}
    \item[\textbf{(ML')}]
    \emph{Para todo índice $\alpha$, existe un $\beta\geq\alpha$ tal
    que, para todo $\gamma\geq\beta$,
    $f_{\alpha\gamma}(A_\gamma)$ es denso en
    $f_{\alpha\beta}(A_\beta)$.}
  \end{itemize}

  Esto da un suplemento análogo a \sref{0.13.2.3}.
  Supongamos que $\mathrm{K}^n_\alpha=0$ para $n<0$ y todo $\alpha$;
  que
  $(\mathrm{K}^n_\alpha)_{\alpha\in\mathrm{I}}$ satisface (ML) para
  $n\geq0$; y que
  $A_\alpha=\HH^0(\mathrm{K}^\bullet_\alpha)$ puede dotarse de
  estructuras de espacio métrico que tengan las propiedades recién enunciadas.
  Entonces las conclusiones de \sref{0.13.2.3} no cambian para
  $n\geq2$, y $h_1$ es además \emph{biyectivo}.
  En efecto, el argumento de \sref{0.13.2.2} muestra también que
  $(\mathrm{B}^1(\mathrm{K}^\bullet_\alpha))_{\alpha\in\mathrm{I}}$
  satisface (ML), que la sucesión $(***_1)$ es exacta y, finalmente,
  que
  \[
    \varprojlim_\alpha
    \mathrm{B}^1(\mathrm{K}^\bullet_\alpha)
    =
    \mathrm{B}^1(\mathrm{K}^\bullet).
  \]
  Esto demuestra, entre otras cosas, las aserciones de (T, 3.10.2).

  \item[{\rm(ii)}]
  Se pueden introducir los funtores derivados por la izquierda
  $\varinjlim^{(n)}$ del funtor $\varinjlim$, y obtener enunciados más
  completos que los precedentes~\cite{III-28}.
\end{enumerate}
\end{env}

\subsection{Aplicación: cohomología de un límite proyectivo de haces}
\label{subsection:0.13.3}

\begin{proposition}[13.3.1]
\label{0.13.3.1}
\oldpage[0\textsubscript{III}]{68}
Sea $X$ un espacio topológico, sea
$(\sh{F}_k)_{k\in\bb{N}}$ un sistema proyectivo de haces de
grupos abelianos sobre $X$, y pongamos
$\sh{F}=\varprojlim_k\sh{F}_k$.
Supongamos que se cumplen las condiciones siguientes:
\begin{enumerate}
  \item[{\rm(i)}]
  Existe una base $\mathfrak{B}$ de la topología de $X$ tal que,
  para todo $U\in\mathfrak{B}$ y todo $i\geq0$, el sistema
  proyectivo
  $(\HH^i(U,\sh{F}_k))_{k\in\bb{N}}$ satisface (ML).

  \item[{\rm(ii)}]
  Para todo $x\in X$ y todo $i>0$, se tiene
  \[
    \varinjlim_U\bigl(\varprojlim_k
    \HH^i(U,\sh{F}_k)\bigr)=0,
  \]
  donde $U$ recorre los entornos de $x$ pertenecientes a
  $\mathfrak{B}$.

  \item[{\rm(iii)}]
  Los homomorfismos
  $u_{hk}:\sh{F}_k\longrightarrow\sh{F}_h$ $(h\leq k)$ que definen el
  sistema proyectivo $(\sh{F}_k)$ son sobreyectivos.
\end{enumerate}
Bajo estas condiciones, para todo $i>0$, el homomorfismo canónico
\[
  h_i:\HH^i(X,\sh{F})
  \longrightarrow
  \varprojlim_k\HH^i(X,\sh{F}_k)
\]
es sobreyectivo. Si, además, para un valor dado de $i$ el sistema
proyectivo
$(\HH^{i-1}(X,\sh{F}_k))_{k\in\bb{N}}$ satisface (ML), entonces $h_i$ es
biyectivo.
\end{proposition}

\begin{proof}
\emph{a)}
Supongamos primero que los $\sh{F}_k$, así como los núcleos
$\sh{N}_{hk}$ de los $u_{hk}$, son flácidos; mostraremos entonces que
la condición~{\rm(iii)} del enunciado implica
$\HH^i(X,\sh{F})=0$ para $i>0$.
Basta demostrar que, para todo subconjunto abierto $U$ de $X$ y todo
recubrimiento $\mathfrak{U}$ de $U$ por subconjuntos abiertos de $U$, se tiene
\[
  \CHH^i(\mathfrak{U},\sh{F})=0
  \qquad (i>0).
\]
En efecto, se sigue entonces primero que, para la cohomología de \v{C}ech,
$\check{\HH}^{\,i}(U,\sh{F})=0$ para todo $i>0$, y por tanto, por
(G, II, 5.9.2) aplicado al conjunto de todos los subconjuntos abiertos de
$X$, que $\HH^i(X,\sh{F})=0$ para todo $i>0$.

Como los $\sh{F}_k$ son flácidos, se tiene
$\CHH^i(\mathfrak{U},\sh{F}_k)=0$ para todo $i>0$
(G, II, 5.2.3).
Para cada $k$, consideremos el complejo
$\check{\mathrm{C}}^\bullet(\mathfrak{U},\sh{F}_k)$ de cocadenas
alternadas sobre el nervio del recubrimiento $\mathfrak{U}$ (G, II, 5.1);
estos complejos forman manifiestamente un sistema proyectivo de complejos de
grupos abelianos.
Demostremos que todas las aplicaciones
\[
  \check{\mathrm{C}}^\bullet(\mathfrak{U},\sh{F}_k)
  \longrightarrow
  \check{\mathrm{C}}^\bullet(\mathfrak{U},\sh{F}_h)
  \qquad(h\leq k)
\]
son sobreyectivas.
Por definición, basta mostrar que, para todo subconjunto abierto $V$
de $X$, la aplicación
$\Gamma(V,\sh{F}_k)\longrightarrow\Gamma(V,\sh{F}_h)$ es
es sobreyectiva. Pero la sucesión
\[
  0\longrightarrow\sh{N}_{hk}\longrightarrow\sh{F}_k
  \longrightarrow\sh{F}_h\longrightarrow0
\]
es exacta por hipótesis y da, por tanto, la sucesión exacta de
cohomología
\[
  \Gamma(V,\sh{F}_k)\longrightarrow\Gamma(V,\sh{F}_h)
  \longrightarrow\HH^1(V,\sh{N}_{hk})=0,
\]
pues $\sh{N}_{hk}$ es flácido.

El sistema proyectivo
$(\check{\mathrm{C}}^\bullet(\mathfrak{U},\sh{F}_k))_{k\in\bb{N}}$
satisface, por tanto, (ML), y lo mismo vale para
$(\CHH^i(\mathfrak{U},\sh{F}_k))_{k\in\bb{N}}$ para todo $i\geq0$:
para $i>0$ esto es trivial, mientras que para $i=0$ se sigue de lo
precedente, pues
$\CHH^0(\mathfrak{U},\sh{F}_k)=\Gamma(U,\sh{F}_k)$
(G, II, 5.2.2).
Podemos, por tanto, aplicar \sref{0.13.2.3}, que da
\[
  \CHH^i(\mathfrak{U},\sh{F})
  =
  \varprojlim_k\CHH^i(\mathfrak{U},\sh{F}_k)
  =0
  \qquad(i>0).
\]

\emph{b)}
Pasemos ahora al caso general. Para cada $k\in\bb{N}$, consideremos la
\emph{resolución canónica}
\[
  \sh{C}^\bullet(X,\sh{F}_k)
  =
  \bigl(\sh{C}^i(X,\sh{F}_k)\bigr)_{i\geq0}
\]
de $\sh{F}_k$ por haces flácidos (G, II, 4.3).
Para todo $i\geq0$, está claro que
$(\sh{C}^i(X,\sh{F}_k))_{k\in\bb{N}}$ es un sistema proyectivo de
haces flácidos; demostremos que satisface las condiciones de la
parte~\emph{a)}.
\oldpage[0\textsubscript{III}]{69}
En efecto, si $\sh{N}_{hk}$ denota el núcleo de $u_{hk}$ para $h\leq k$,
entonces la sucesión
\[
  0\longrightarrow\sh{N}_{hk}\longrightarrow\sh{F}_k
  \longrightarrow\sh{F}_h\longrightarrow0
\]
es exacta por~{\rm(iii)}, y nuestra aserción se sigue de la exactitud
del funtor
$\sh{A}\longmapsto\sh{C}^i(X,\sh{A})$ (G, II, 4.3).

Pongamos
\[
  \sh{G}^i=\varprojlim_k\sh{C}^i(X,\sh{F}_k).
\]
Por la parte~\emph{a)}, se tiene
$\HH^j(X,\sh{G}^i)=0$ para $j>0$ e $i\geq0$.
Demostraremos que
$\sh{G}^\bullet=(\sh{G}^i)_{i\geq0}$ es una resolución del haz
$\sh{F}$. Como $\HH^j(X,\sh{G}^i)=0$ para $j>0$, la cohomología
$\HH^\bullet(X,\sh{F})$ será entonces igual a
$\HH^\bullet(\Gamma(X,\sh{G}^\bullet))$ (G, II, 4.7.1).

Tomando el límite proyectivo de las sucesiones exactas
\[
  0\longrightarrow\sh{F}_k
  \longrightarrow\sh{C}^0(X,\sh{F}_k)
  \longrightarrow\sh{C}^1(X,\sh{F}_k)
  \longrightarrow\cdots
\]
se obtiene manifiestamente un complejo de haces de grupos abelianos
\[
  0\longrightarrow\sh{F}
  \longrightarrow\sh{G}^0
  \longrightarrow\sh{G}^1
  \longrightarrow\cdots.
\]
Para demostrar nuestra aserción, queda establecer que
$\sh{H}^i(\sh{G}^\bullet)=0$ para $i>0$.
Este haz está generado por el prehaz
\[
  U\longmapsto\HH^i\bigl(\Gamma(U,\sh{G}^\bullet)\bigr)
\]
(G, II, 4.1).
Ahora bien, el complejo $\Gamma(U,\sh{G}^\bullet)$ es el límite proyectivo
del sistema proyectivo de complejos de grupos abelianos
\[
  \bigl(\Gamma(U,\sh{C}^\bullet(X,\sh{F}_k))\bigr)_{k\in\bb{N}}
\]
(0\textsubscript{I}, \EGAUnavailableHyperref{0.3.2.6}{3.2.6}).
Vimos en la parte~\emph{a)} que, para todo $i\geq0$, las aplicaciones
\[
  \Gamma(U,\sh{C}^i(X,\sh{F}_k))
  \longrightarrow
  \Gamma(U,\sh{C}^i(X,\sh{F}_h))
  \qquad(h\leq k)
\]
son sobreyectivas. Por otra parte,
\[
  \HH^i(U,\sh{F}_k)
  =
  \HH^i\bigl(\Gamma(U,\sh{C}^\bullet(X,\sh{F}_k))\bigr),
\]
pues la resolución canónica $\sh{C}^\bullet(U,\sh{F}_k|U)$ es
inducida sobre $U$ por $\sh{C}^\bullet(X,\sh{F}_k)$.
Por la hipótesis~{\rm(i)}, para todo $U\in\mathfrak{B}$ podemos aplicar
\sref{0.13.2.3} a este sistema proyectivo de complejos, y por tanto
\[
  \HH^i\bigl(\Gamma(U,\sh{G}^\bullet)\bigr)
  =
  \varprojlim_k\HH^i(U,\sh{F}_k)
  \qquad(i\geq0).
\]
La hipótesis~{\rm(ii)} muestra entonces, por definición, que los haces
$\sh{H}^i(\sh{G}^\bullet)$ se anulan para $i>0$.

Así, para todo $i\geq0$,
\[
  \HH^i(X,\sh{F})
  =
  \HH^i\bigl(\Gamma(X,\sh{G}^\bullet)\bigr),
  \qquad
  \Gamma(X,\sh{G}^\bullet)
  =
  \varprojlim_k\Gamma(X,\sh{C}^\bullet(X,\sh{F}_k)).
\]
Acabamos de observar que todas las aplicaciones
\[
  \Gamma(X,\sh{C}^i(X,\sh{F}_k))
  \longrightarrow
  \Gamma(X,\sh{C}^i(X,\sh{F}_h))
  \qquad(h\leq k)
\]
son sobreyectivas; la conclusión se sigue, por tanto, una vez más de
\sref{0.13.2.3}.
\end{proof}

\begin{env}[13.3.2]
\label{0.13.3.2}
\emph{Observaciones} (13.3.2).
\begin{enumerate}
  \item[{\rm(i)}]
  El enunciado de \sref{0.13.3.1} solo tiene interés para $i>0$,
  pues, para $i=0$, $h_i$ es siempre un isomorfismo sin hipótesis
  alguna (0\textsubscript{I}, \EGAUnavailableHyperref{0.3.2.6}{3.2.6}).

  \item[{\rm(ii)}]
  Las condiciones~{\rm(i)} y~{\rm(ii)} de \sref{0.13.3.1} se satisfacen,
  en particular, si
  $\HH^i(U,\sh{F}_k)=0$ para todo $k$, todo $i>0$ y todo
  $U\in\mathfrak{B}$, y si, para $U\in\mathfrak{B}$, las aplicaciones
  \[
    \Gamma(U,\sh{F}_k)\longrightarrow\Gamma(U,\sh{F}_h)
  \]
  son sobreyectivas. Este es el caso más frecuente en el que se aplica
  \sref{0.13.3.1}.
\end{enumerate}
\end{env}

\subsection{Condición de Mittag-Leffler y objetos graduados asociados a sistemas proyectivos}
\label{subsection:0.13.4}
\label{0.13.4}

\begin{env}[13.4.1]
\label{0.13.4.1}
Sea
$\bb{A}=(A_k,u_{kh})_{k\in\bb{Z}}$ un sistema proyectivo en una
categoría abeliana $\cat{C}$.
Decimos que está \emph{acotado inferiormente} si existe un $k_0$ tal que
$A_k=0$ para $k<k_0$.
Sobre cada $A_k$ definimos una \emph{filtración}
$(\mathrm{F}^p(A_k))_{p\in\bb{Z}}$ por
\begin{equation}
\label{0.13.4.1.1}
\tag{13.4.1.1}
  \mathrm{F}^p(A_k)=
  \begin{cases}
    \Ker(A_k\longrightarrow A_{p-1}) & \text{para }p\leq k+1,\\
    0 & \text{para }p\geq k+1.
  \end{cases}
\end{equation}

\oldpage[0\textsubscript{III}]{70}
Por hipótesis,
$\mathrm{F}^{k_0}(A_k)=A_k$ y
$\mathrm{F}^{k+1}(A_k)=0$; en otras palabras, la filtración considerada
es finita (\sref{0.11.1.3}).
Los objetos graduados asociados a esta filtración son, por tanto,
\[
  \mathrm{gr}^p(A_k)
  =
  \Ker(A_k\longrightarrow A_{p-1})
  \big/
  \Ker(A_k\longrightarrow A_p).
\]
Por consiguiente, $\mathrm{gr}^p(A_k)$ es isomorfo a la imagen por
$A_k\longrightarrow A_p$ de
$\Ker(A_k\longrightarrow A_{p-1})$.
Por la transitividad de los morfismos que definen un sistema proyectivo,
se tiene, por tanto,
\begin{equation}
\label{0.13.4.1.2}
\tag{13.4.1.2}
  \mathrm{gr}^p(A_k)
  =
  \Ker(A_p\longrightarrow A_{p-1})
  \cap
  \Im(A_k\longrightarrow A_p).
\end{equation}
Pero \sref{0.13.4.1.1} da
$\Ker(A_p\longrightarrow A_{p-1})=\mathrm{gr}^p(A_p)$, y por tanto
\begin{equation}
\label{0.13.4.1.3}
\tag{13.4.1.3}
  \mathrm{gr}^p(A_k)
  =
  \mathrm{gr}^p(A_p)
  \cap
  \Im(A_k\longrightarrow A_p).
\end{equation}

Las definiciones precedentes muestran también que, para $k\leq h$,
\[
  u_{kh}\bigl(\mathrm{F}^p(A_h)\bigr)
  \subset
  \mathrm{F}^p(A_k),
\]
y, por consiguiente, para todo $p\in\bb{Z}$, los
$\mathrm{gr}^p(u_{kh})$ definen un sistema proyectivo
$(\mathrm{gr}^p(A_k))_{k\in\bb{Z}}$.
\end{env}

\begin{env}[13.4.2]
\label{0.13.4.2}
Decimos que el sistema proyectivo $\bb{A}$ es
\emph{esencialmente constante} si los morfismos
$A_{k+1}\longrightarrow A_k$ son isomorfismos para todo $k$
suficientemente grande.
Decimos que $\bb{A}$ es \emph{estricto} si los morfismos
$A_i\longrightarrow A_j$ $(j\leq i)$ son epimorfismos.

Cuando $\bb{A}$ es estricto, se sigue de \sref{0.13.4.1.3} que, para
$p\leq k\leq h$, el morfismo canónico
\[
  \mathrm{gr}^p(A_h)\longrightarrow\mathrm{gr}^p(A_k)
\]
es un isomorfismo; en otras palabras, el sistema proyectivo
$(\mathrm{gr}^p(A_k))_{k\in\bb{Z}}$ es esencialmente constante.
La sucesión de objetos $\mathrm{gr}^p(A_p)$, identificada para todo
$p$ con $\varprojlim_k\mathrm{gr}^p(A_k)$, se denota entonces
$\mathrm{gr}^\bullet(\bb{A})$ y se llama el
\emph{objeto graduado asociado al sistema proyectivo estricto}
$\bb{A}=(A_k)$.

Supongamos ahora que el sistema proyectivo $\bb{A}$, acotado
inferiormente, satisface (ML).
Por \sref{0.13.1.2}, el sistema proyectivo
$\bb{A}'=(A'_k)$ de objetos de ``imagen universal'' es estricto; también está
acotado inferiormente.
El objeto graduado $\mathrm{gr}^\bullet(\bb{A}')$ asociado a
$\bb{A}'$ se llama asimismo el objeto graduado
\emph{asociado a} $\bb{A}$ y se denota
$\mathrm{gr}^\bullet(\bb{A})$.
\end{env}

\begin{proposition}[13.4.3]
\label{0.13.4.3}
Sea $\bb{A}=(A_k)_{k\in\bb{Z}}$ un sistema proyectivo acotado
inferiormente en una categoría abeliana $\cat{C}$.
Las condiciones siguientes son equivalentes:
\begin{enumerate}
  \item[{\rm a)}] $\bb{A}$ satisface (ML).
  \item[{\rm b)}] Para todo $p\in\bb{Z}$, el sistema proyectivo
  $(\mathrm{gr}^p(A_k))_{k\in\bb{Z}}$ es esencialmente constante.
\end{enumerate}
Además, cuando se cumplen estas condiciones, existe, para todo
$p\in\bb{Z}$, un isomorfismo canónico
\begin{equation}
\label{0.13.4.3.1}
\tag{13.4.3.1}
  \mathrm{gr}^p(\bb{A})
  \isoto
  \varprojlim_k\mathrm{gr}^p(A_k).
\end{equation}
\end{proposition}

\begin{proof}
Se sigue inmediatamente de \sref{0.13.4.1.2} que \emph{a)} implica
\emph{b)}. La misma fórmula, aplicada al sistema proyectivo
$\bb{A}'$ con la notación de \sref{0.13.4.2}, da
\sref{0.13.4.3.1} por definición.

Para $k\leq h$, pongamos
\[
  A_{kh}=\Im(A_h\longrightarrow A_k).
\]
Si $k\leq h\leq j$, entonces
$A_{kj}\subset A_{kh}\subset A_k$.
Dotemos a $A_{kh}$ de la filtración inducida por
$(\mathrm{F}^p(A_k))$.
La transitividad de los morfismos que definen $\bb{A}$ muestra
inmediatamente que esta es también la filtración cociente de
$(\mathrm{F}^p(A_h))$; por consiguiente,
\begin{equation}
\label{0.13.4.3.2}
\tag{13.4.3.2}
  \mathrm{gr}^p(A_{kh})
  =
  \Im\bigl(
    \mathrm{gr}^p(A_h)\longrightarrow\mathrm{gr}^p(A_k)
  \bigr).
\end{equation}

\oldpage[0\textsubscript{III}]{71}
Supongamos ahora que se cumple \emph{b)}.
Para todo $p\in\bb{Z}$ y todo $k\geq p$, existe un entero
$\mathrm{L}(p,k)$ tal que el miembro derecho de
\sref{0.13.4.3.2} es constante para $h\geq\mathrm{L}(p,k)$.
Como $\mathrm{gr}^p(A_k)=0$ para $p<k_0$, para $k$ fijo solo
un número finito de los enteros $\mathrm{L}(p,k)$ puede ser no nulo cuando
$p$ recorre los enteros $\leq k$.
Sea $\mathrm{L}(k)=m$ el mayor de estos enteros.
Para todo $h\geq m$, se tiene
$A_{kh}\subset A_{km}$, y, por la definición de $m$, la inyección
canónica
$A_{kh}\longrightarrow A_{km}$ induce un isomorfismo
\[
  \mathrm{gr}^\bullet(A_{kh})
  \simto
  \mathrm{gr}^\bullet(A_{km}).
\]
Como las filtraciones son finitas, esta inyección es ella misma biyectiva
(Bourbaki, \emph{Alg.\ comm.}, cap.~III, \textsection2,
n\textsuperscript{o}~8, teorema~1).
Esto demuestra que $\bb{A}$ satisface (ML).
\end{proof}

\begin{env}[13.4.4]
\label{0.13.4.4}
Supongamos que el límite proyectivo
$A=\varprojlim_k A_k$ existe en $\cat{C}$.
En las definiciones de \sref{0.13.4.1}, se puede entonces reemplazar $A_k$
por $A$, y la filtración así definida sobre $A$ sigue satisfaciendo
\begin{equation}
\label{0.13.4.4.1}
\tag{13.4.4.1}
  \mathrm{gr}^p(A)
  =
  \mathrm{gr}^p(A_p)
  \cap
  \Im(A\longrightarrow A_p).
\end{equation}
\end{env}

\begin{corollary}[13.4.5]
\label{0.13.4.5}
Supongamos que $\cat{C}$ es la categoría de grupos abelianos.
Si el sistema proyectivo $\bb{A}$ satisface (ML) y
$A=\varprojlim_k A_k$, entonces para todo $p\in\bb{Z}$ existe un
isomorfismo canónico
\begin{equation}
\label{0.13.4.5.1}
\tag{13.4.5.1}
  \mathrm{gr}^p(A)
  \isoto
  \varprojlim_k\mathrm{gr}^p(A_k).
\end{equation}
\end{corollary}

\begin{proof}
En efecto,
\[
  \Im(A_k\longrightarrow A_p)
  =
  \Im(A\longrightarrow A_p)
\]
para todo $k$ suficientemente grande
(Bourbaki, \emph{Top.\ g\'{e}n.}, cap.~II,
3\textsuperscript{e}~\'{e}d., \textsection3,
n\textsuperscript{o}~5, teorema~1), y la conclusión se sigue de
\sref{0.13.4.1.3} y \sref{0.13.4.4.1}.
\end{proof}

\subsection{Límites proyectivos de sucesiones espectrales de complejos filtrados}
\label{subsection:0.13.5}

\begin{env}[13.5.1]
\label{0.13.5.1}
Sea $\cat{C}$ una categoría abeliana y sea $X^\bullet$ un complejo
de objetos de $\cat{C}$ dotado de una filtración
$(\mathrm{F}^p(X^\bullet))_{p\in\bb{Z}}$ tal que
$\mathrm{F}^{p_0}(X^\bullet)=X^\bullet$ para algún índice $p_0$.
Para todo $k\in\bb{Z}$, consideremos el complejo
\[
  X^\bullet_k=X^\bullet/\mathrm{F}^{k+1}(X^\bullet);
\]
está dotado canónicamente de la filtración constituida por
\[
  \mathrm{F}^p(X^\bullet_k)
  =
  \mathrm{F}^p(X^\bullet)/\mathrm{F}^{k+1}(X^\bullet)
  \quad\text{para }p\leq k,
\]
y $\mathrm{F}^p(X^\bullet_k)=0$ para $p\geq k+1$.
Además, existen morfismos canónicos
$X^\bullet_{k+1}\longrightarrow X^\bullet_k$, que hacen de
$\bb{X}^\bullet=(X^\bullet_k)_{k\in\bb{Z}}$ un \emph{sistema proyectivo
de complejos filtrados de objetos de $\cat{C}$}.
Observemos que este sistema proyectivo es \emph{estricto} y que
$X^\bullet_k=0$ para $k<p_0$.
\end{env}

\begin{env}[13.5.2]
\label{0.13.5.2}
Más generalmente, consideremos un sistema proyectivo estricto acotado
inferiormente
$\bb{X}^\bullet=(X^\bullet_k)_{k\in\bb{Z}}$ de complejos de objetos de
$\cat{C}$, y dotemos cada $X^\bullet_k$ de la filtración definida en
\sref{0.13.4.1} (aplicada en la categoría abeliana de complejos acotados
inferiormente de objetos de $\cat{C}$).
Los morfismos
$X^\bullet_k\longrightarrow X^\bullet_p$ ($p\leq k$) se convierten entonces
en morfismos de complejos \emph{filtrados} con filtraciones finitas.
La funtorialidad de las sucesiones espectrales de complejos filtrados
(\sref{0.11.2.3}) muestra que los morfismos de transición de
$\bb{X}^\bullet$ inducen morfismos que hacen de
\[
  \mathrm{E}(\bb{X}^\bullet)
  =
  \bigl(\mathrm{E}(X^\bullet_k)\bigr)_{k\in\bb{Z}}
\]
un \emph{sistema proyectivo de sucesiones espectrales}.
\end{env}

\begin{lemma}[13.5.3]
\label{0.13.5.3}
Supongamos que el sistema proyectivo
$\bb{X}^\bullet=(X^\bullet_k)_{k\in\bb{Z}}$ de complejos filtrados se
obtiene como en \sref{0.13.5.2}. Entonces:
\begin{enumerate}
  \item[{\rm a)}] Para $r\geq p-p_0$, se tiene
  \[
    \mathrm{B}_r^{pq}(X^\bullet_k)
    =
    \mathrm{B}_\infty^{pq}(X^\bullet_k)
  \]
  para todo $k\in\bb{Z}$.
  \item[{\rm b)}] Para $k+1\leq p+r$, se tiene
  \[
    \mathrm{Z}_r^{pq}(X^\bullet_k)
    =
    \mathrm{Z}_\infty^{pq}(X^\bullet_k).
  \]
  \item[{\rm c)}] Para $k+1\geq p+r$, los morfismos
  \[
    \mathrm{Z}_r^{pq}(X^\bullet_h)
    \longrightarrow
    \mathrm{Z}_r^{pq}(X^\bullet_k)
    \quad\text{y}\quad
    \mathrm{B}_r^{pq}(X^\bullet_h)
    \longrightarrow
    \mathrm{B}_r^{pq}(X^\bullet_k)
  \]
  son isomorfismos para todo $h\geq k$.
\end{enumerate}
\end{lemma}

\oldpage[0\textsubscript{III}]{72}
Estas tres propiedades se siguen inmediatamente de las definiciones de
\sref{0.11.2.2}, teniendo en cuenta que
$\mathrm{F}^{p-r+1}(X^\bullet_k)=X^\bullet_k$ cuando $p-r<p_0$.

\begin{env}[13.5.4]
\label{0.13.5.4}
Supongamos que se cumplen las hipótesis de \sref{0.13.5.3}.
Entonces, para $p$, $q$ y $r$ fijos (con $r$ finito), los sistemas
proyectivos
\[
  \bigl(\mathrm{Z}_r^{pq}(X^\bullet_k)\bigr)_{k\in\bb{Z}},
  \qquad
  \bigl(\mathrm{B}_r^{pq}(X^\bullet_k)\bigr)_{k\in\bb{Z}},
  \qquad
  \bigl(\mathrm{E}_r^{pq}(X^\bullet_k)\bigr)_{k\in\bb{Z}}
\]
son \emph{esencialmente constantes}; sus respectivos límites proyectivos
se denotarán por
$\mathrm{Z}_r^{pq}(\bb{X}^\bullet)$,
$\mathrm{B}_r^{pq}(\bb{X}^\bullet)$, y
\[
  \mathrm{E}_r^{pq}(\bb{X}^\bullet)
  =
  \mathrm{Z}_r^{pq}(\bb{X}^\bullet)/
  \mathrm{B}_r^{pq}(\bb{X}^\bullet).
\]
Los objetos $\mathrm{Z}_r^{pq}(\bb{X}^\bullet)$ y
$\mathrm{B}_r^{pq}(\bb{X}^\bullet)$ se identifican canónicamente con
subobjetos de $\mathrm{E}_1^{pq}(\bb{X}^\bullet)$.
La definición de los $d_r^{pq}$ (M, XV, 1) muestra que estos
morfismos, tomados respecto de los $X^\bullet_k$, son también esencialmente
constantes, y definen, por tanto, morfismos
\begin{equation}
\label{0.13.5.4.1}
\tag{13.5.4.1}
  d_r^{pq}:
  \mathrm{E}_r^{pq}(\bb{X}^\bullet)
  \longrightarrow
  \mathrm{E}_r^{p+r,q-r+1}(\bb{X}^\bullet)
\end{equation}
tal que
$d_r^{p+r,q-r+1}\circ d_r^{pq}=0$.
Además, existen isomorfismos canónicos de
$\Ker(d_r^{pq})$ sobre
\[
  \mathrm{Z}_{r+1}^{pq}(\bb{X}^\bullet)/
  \mathrm{B}_r^{pq}(\bb{X}^\bullet)
\]
y de $\Im(d_r^{pq})$ sobre
\[
  \mathrm{B}_{r+1}^{p+r,q-r+1}(\bb{X}^\bullet)/
  \mathrm{B}_r^{p+r,q-r+1}(\bb{X}^\bullet).
\]
\end{env}

\begin{lemma}[13.5.5]
\label{0.13.5.5}
Bajo las hipótesis de \sref{0.13.5.3}, para
$s\geq r>p-p_0$ existe un monomorfismo canónico
\begin{equation}
\label{0.13.5.5.1}
\tag{13.5.5.1}
  i:
  \mathrm{E}_s^{pq}(\bb{X}^\bullet)
  \longrightarrow
  \mathrm{E}_r^{pq}(\bb{X}^\bullet)
\end{equation}
y un isomorfismo canónico
\begin{equation}
\label{0.13.5.5.2}
\tag{13.5.5.2}
  j_r:
  \mathrm{E}_r^{pq}(\bb{X}^\bullet)
  \simto
  \mathrm{E}_\infty^{pq}(X^\bullet_{p+r-1})
\end{equation}
tal que el diagrama
\begin{equation}
\label{0.13.5.5.3}
\tag{13.5.5.3}
  \xymatrix{
    \mathrm{E}_s^{pq}(\bb{X}^\bullet)
      \ar[r]^{j_s}\ar[d]_i
    &
    \mathrm{E}_\infty^{pq}(X^\bullet_{p+s-1})
      \ar[d]
    \\
    \mathrm{E}_r^{pq}(\bb{X}^\bullet)
      \ar[r]^{j_r}
    &
    \mathrm{E}_\infty^{pq}(X^\bullet_{p+r-1})
  }
\end{equation}
es conmutativo, donde la flecha vertical derecha está inducida por el
morfismo
$X^\bullet_{p+s-1}\longrightarrow X^\bullet_{p+r-1}$.
\end{lemma}

\begin{proof}
La existencia de $i$ se sigue de
\[
  \mathrm{B}_r^{pq}(X^\bullet_k)
  =
  \mathrm{B}_\infty^{pq}(X^\bullet_k)
  \qquad (r>p-p_0)
\]
(\sref{0.13.5.3},~a)).
Para $k+1\leq p+r$, se tiene
\[
  \mathrm{Z}_r^{pq}(X^\bullet_k)
  =
  \mathrm{Z}_\infty^{pq}(X^\bullet_k)
\]
(\sref{0.13.5.3},~b)); en particular,
\[
  \mathrm{Z}_r^{pq}(X^\bullet_{p+r-1})
  =
  \mathrm{Z}_\infty^{pq}(X^\bullet_{p+r-1}).
\]
Por otra parte,
$\mathrm{Z}_r^{pq}(X^\bullet_{p+r-1})$ y
$\mathrm{B}_r^{pq}(X^\bullet_{p+r-1})$ se identifican canónicamente con
$\mathrm{Z}_r^{pq}(\bb{X}^\bullet)$ y
$\mathrm{B}_r^{pq}(\bb{X}^\bullet)$ por
\sref{0.13.5.3},~c).
Esto demuestra la existencia de $j_r$ y la conmutatividad de
\sref{0.13.5.5.1}.
\end{proof}

\begin{corollary}[13.5.6]
\label{0.13.5.6}
Bajo las hipótesis de \sref{0.13.5.3}, si uno de los límites
proyectivos
\[
  \varprojlim_r\mathrm{E}_r^{pq}(\bb{X}^\bullet),
  \qquad
  \varprojlim_k\mathrm{E}_\infty^{pq}(X^\bullet_k)
\]
existe, entonces también existe el otro, y hay un isomorfismo canónico
\begin{equation}
\label{0.13.5.6.1}
\tag{13.5.6.1}
  j_\infty:
  \varprojlim_r\mathrm{E}_r^{pq}(\bb{X}^\bullet)
  \simto
  \varprojlim_k\mathrm{E}_\infty^{pq}(X^\bullet_k).
\end{equation}
Además, el sistema proyectivo
$(\mathrm{E}_r^{pq}(\bb{X}^\bullet))_{r\in\bb{Z}}$ es esencialmente
constante (\sref{0.13.4.2}) si y solo si el sistema proyectivo
$(\mathrm{E}_\infty^{pq}(X^\bullet_k))_{k\in\bb{Z}}$ es esencialmente
constante.
\end{corollary}

\begin{env}[13.5.7]
\label{0.13.5.7}
Denotamos por
$\mathrm{B}_\infty^{pq}(\bb{X}^\bullet)$ y
$\mathrm{Z}_\infty^{pq}(\bb{X}^\bullet)$ los subobjetos de
$\mathrm{E}_1^{pq}(\bb{X}^\bullet)$ iguales respectivamente a
$\mathrm{B}_r^{pq}(\bb{X}^\bullet)$ para $r>p-p_0$ y a
$\inf_r\mathrm{Z}_r^{pq}(\bb{X}^\bullet)$, cuando este último existe.
Así,
\[
  \varprojlim_r\mathrm{E}_r^{pq}(\bb{X}^\bullet)
\]
\oldpage[0\textsubscript{III}]{73}
se identifica canónicamente con
\[
  \mathrm{E}_\infty^{pq}(\bb{X}^\bullet)
  =
  \mathrm{Z}_\infty^{pq}(\bb{X}^\bullet)/
  \mathrm{B}_\infty^{pq}(\bb{X}^\bullet).
\]
Observemos que los objetos
$\mathrm{Z}_r^{pq}(\bb{X}^\bullet)$,
$\mathrm{B}_r^{pq}(\bb{X}^\bullet)$,
$\mathrm{E}_r^{pq}(\bb{X}^\bullet)$
($1\leq r\leq+\infty$), y los morfismos $d_r^{pq}$ dependen
\emph{funtorialmente} del sistema proyectivo $\bb{X}^\bullet$, con las
restricciones de \sref{0.13.5.5}, y los morfismos definidos en
\sref{0.13.5.5} y \sref{0.13.5.6} son funtoriales.
\end{env}

\subsection{Sucesión espectral de un funtor relativa a un objeto dotado de una filtración finita}
\label{subsection:0.13.6}

\begin{env}[13.6.1]
\label{0.13.6.1}
Sean $\cat{C}$ y $\cat{C}'$ dos categorías abelianas, y sea
$\mathrm{T}:\cat{C}\longrightarrow\cat{C}'$ un funtor covariante
aditivo.
Supongamos que todo objeto de $\cat{C}$ es isomorfo a un subobjeto de
un objeto inyectivo, de modo que existan los funtores derivados por la derecha
$\RR^p\mathrm{T}$ ($p\geq0$).
\end{env}

\begin{lemma}[13.6.2]
\label{0.13.6.2}
Sea $A$ un objeto de $\cat{C}$ dotado de una filtración finita
$(\mathrm{F}^i(A))_{i\in\bb{Z}}$.
Existe una resolución inyectiva
$X^\bullet=(X^j)_{j\geq0}$ de $A$, dotada de una filtración finita
$(\mathrm{F}^i(X^\bullet))_{i\in\bb{Z}}$, tal que
\[
  \mathrm{F}^i(A)=A
  \quad\text{(respectivamente, }\mathrm{F}^i(A)=0\text{)}
\]
implica
\[
  \mathrm{F}^i(X^\bullet)=X^\bullet
  \quad\text{(respectivamente, }\mathrm{F}^i(X^\bullet)=0\text{)},
\]
y tal que, para todo $i\in\bb{Z}$,
$\mathrm{F}^i(X^\bullet)$ es una resolución inyectiva de
$\mathrm{F}^i(A)$.
\end{lemma}

\begin{proof}
Sea $p$ (respectivamente, $q>p$) el mayor índice tal que
$\mathrm{F}^p(A)=A$ (respectivamente, el menor índice tal que
$\mathrm{F}^q(A)=0$).
Razonamos por inducción sobre $q-p$, siendo inmediato el lema cuando
$q-p=1$.
Supongamos construida una resolución inyectiva ${X''}^\bullet$ de
$A/\mathrm{F}^{q-1}(A)$ que tenga las propiedades requeridas.
Consideremos la sucesión exacta
\[
  0\longrightarrow\mathrm{F}^{q-1}(A)
  \longrightarrow A
  \longrightarrow A/\mathrm{F}^{q-1}(A)
  \longrightarrow0,
\]
tomemos una resolución inyectiva ${X'}^\bullet$ de
$\mathrm{F}^{q-1}(A)$, y elijamos después una resolución inyectiva
$X^\bullet$ de $A$ de modo que se obtenga una sucesión exacta
\[
  0\longrightarrow {X'}^\bullet
  \longrightarrow X^\bullet
  \longrightarrow {X''}^\bullet
  \longrightarrow0
\]
compatible con la sucesión precedente (M, V, 2.2).
Está claro que $X^\bullet$ tiene las propiedades requeridas.
\end{proof}

\begin{corollary}[13.6.3]
\label{0.13.6.3}
Sea $B$ un segundo objeto de $\cat{C}$, dotado de una filtración
finita $(\mathrm{F}^i(B))_{i\in\bb{Z}}$, sea $s$ un entero,
y sea $u:A\longrightarrow B$ un morfismo tal que
\[
  u(\mathrm{F}^i(A))\subset\mathrm{F}^{i+s}(B)
\]
para todo $i\in\bb{Z}$.
Si $Y^\bullet=(Y^j)_{j\geq0}$ es una resolución inyectiva de $B$
dotada de una filtración
$(\mathrm{F}^i(Y^\bullet))_{i\in\bb{Z}}$ que tenga las propiedades
enunciadas en \sref{0.13.6.2}, entonces existe un morfismo
$v:X^\bullet\longrightarrow Y^\bullet$, compatible con $u$, tal
que
\[
  v(\mathrm{F}^i(X^\bullet))
  \subset
  \mathrm{F}^{i+s}(Y^\bullet)
\]
para todo $i\in\bb{Z}$.
Además, cualesquiera dos de tales morfismos $v$ y $v'$ son homotópicos.
\end{corollary}

Esto se sigue inmediatamente, por inducción sobre $q-p$, de la construcción
precedente y de (M, V, 2.3).

\begin{env}[13.6.4]
\label{0.13.6.4}
Bajo las hipótesis de \sref{0.13.6.2}, consideremos el complejo
$\mathrm{T}(X^\bullet)$ en $\cat{C}'$.
Está filtrado por los complejos
$\mathrm{T}(\mathrm{F}^i(X^\bullet))$, pues
$\mathrm{F}^i(X^\bullet)$ es un sumando directo de $X^\bullet$.
Se sigue de \sref{0.13.6.3} que la \emph{sucesión espectral} de
este complejo filtrado depende, salvo isomorfismo, únicamente del
objeto filtrado $A$.
Su término límite es la cohomología $\RR^\bullet\mathrm{T}(A)$, dotada
de la filtración
\begin{equation}
\label{0.13.6.4.1}
\tag{13.6.4.1}
  \begin{split}
  \mathrm{F}^p(\RR^n\mathrm{T}(A))
  &=
  \Im\bigl(
    \RR^n\mathrm{T}(\mathrm{F}^p(A))
    \longrightarrow
    \RR^n\mathrm{T}(A)
  \bigr)
  \\
  &=
  \Ker\bigl(
    \RR^n\mathrm{T}(A)
    \longrightarrow
    \RR^n\mathrm{T}(A/\mathrm{F}^p(A))
  \bigr),
  \end{split}
\end{equation}
por \sref{0.11.2.2}, y su término $\mathrm{E}_1$ es
\begin{equation}
\label{0.13.6.4.2}
\tag{13.6.4.2}
  \mathrm{E}_1^{pq}
  =
  \RR^{p+q}\mathrm{T}(\mathrm{gr}^p(A)),
\end{equation}
donde, como de costumbre,
$\mathrm{gr}^p(A)=\mathrm{F}^p(A)/\mathrm{F}^{p+1}(A)$.
Se desprende de \sref{0.11.2.2} que la filtración del término límite
es finita y que, para $p$ y $q$ fijos, las sucesiones
\oldpage[0\textsubscript{III}]{74}
\[
  \mathrm{B}_r^{pq}(A)
  =
  \mathrm{B}_r^{pq}(\mathrm{T}(X^\bullet))
  \quad\text{y}\quad
  \mathrm{Z}_r^{pq}(A)
  =
  \mathrm{Z}_r^{pq}(\mathrm{T}(X^\bullet))
\]
son estacionarias.
La sucesión espectral precedente es, por tanto, \emph{birregular}
(\sref{0.11.1.3}).
La denotamos por
$\mathrm{E}(A)=(\mathrm{E}_r^{pq}(A))$ y la llamamos la
\emph{sucesión espectral del funtor $\mathrm{T}$ relativa al
objeto filtrado $A$}.
\end{env}

\begin{env}[13.6.5]
\label{0.13.6.5}
Supongamos ahora que se cumplen las hipótesis de \sref{0.13.6.3}, y
conservemos su notación.
Como $\mathrm{F}^i(X^\bullet)$ (respectivamente,
$\mathrm{F}^i(Y^\bullet)$) es un sumando directo de $X^\bullet$
(respectivamente, $Y^\bullet$), se tiene
\[
  (\mathrm{T}v)\bigl(\mathrm{T}(\mathrm{F}^i(X^\bullet))\bigr)
  \subset
  \mathrm{T}(\mathrm{F}^{i+s}(Y^\bullet))
\]
para todo $i\in\bb{Z}$.
Las definiciones de \sref{0.11.2.2} muestran entonces que, para
$1\leq r\leq+\infty$, $\mathrm{T}v$ define morfismos
\[
  \mathrm{B}_r^{pq}(\mathrm{T}(X^\bullet))
  \longrightarrow
  \mathrm{B}_r^{p+s,q-s}(\mathrm{T}(Y^\bullet))
\]
y
\[
  \mathrm{Z}_r^{pq}(\mathrm{T}(X^\bullet))
  \longrightarrow
  \mathrm{Z}_r^{p+s,q-s}(\mathrm{T}(Y^\bullet)),
\]
y, por tanto, un morfismo
\[
  w_r:
  \mathrm{E}_r^{pq}(A)
  \longrightarrow
  \mathrm{E}_r^{p+s,q-s}(B).
\]
Análogamente, sobre los términos límite se tienen morfismos
\[
  u_n:
  \RR^n\mathrm{T}(A)
  \longrightarrow
  \RR^n\mathrm{T}(B)
\]
tales que
\[
  u_n\bigl(\mathrm{F}^p(\RR^n\mathrm{T}(A))\bigr)
  \subset
  \mathrm{F}^{p+s}(\RR^n\mathrm{T}(B)).
\]
La definición de los $d_r^{pq}$ (M, XV, 1) muestra también que los
diagramas
\[
  \xymatrix{
    \mathrm{E}_r^{pq}(A)
      \ar[r]^{d_r^{pq}}\ar[d]_{w_r}
    &
    \mathrm{E}_r^{p+r,q-r+1}(A)
      \ar[d]^{w_r}
    \\
    \mathrm{E}_r^{p+s,q-s}(B)
      \ar[r]_{d_r^{p+s,q-s}}
    &
    \mathrm{E}_r^{p+r+s,q-r-s+1}(B)
  }
\]
son conmutativos.
Se sigue que existe un diagrama conmutativo análogo para los
isomorfismos $\alpha_r^{pq}$, cuya explicitación dejamos al lector.
Finalmente (\emph{loc.\ cit.}), se tienen también diagramas conmutativos
sobre los términos límite:
\[
  \xymatrix{
    \mathrm{E}_\infty^{pq}(A)
      \ar[r]^{\beta^{pq}}\ar[d]_{w_\infty}
    &
    \mathrm{gr}^p(\RR^{p+q}\mathrm{T}(A))
      \ar[d]^{u_{p+q}}
    \\
    \mathrm{E}_\infty^{p+s,q-s}(B)
      \ar[r]_{\beta^{p+s,q-s}}
    &
    \mathrm{gr}^{p+s}(\RR^{p+q}\mathrm{T}(B)).
  }
\]
\end{env}

\begin{env}[13.6.6]
\label{0.13.6.6}
Supongamos, en particular, que existe un anillo $S$, dotado de una
\emph{filtración} $(\mathrm{F}^i(S))_{i\in\bb{Z}}$, y un homomorfismo
de anillos
\begin{equation}
\label{0.13.6.6.1}
\tag{13.6.6.1}
  h:S\longrightarrow\Hom_{\cat{C}}(A,A)
\end{equation}
\oldpage[0\textsubscript{III}]{75}
tal que, para todo $t\in\mathrm{F}^i(S)$,
\[
  h_t(\mathrm{F}^j(A))
  \subset
  \mathrm{F}^{i+j}(A)
\]
para todo par $i,j$.
Para abreviar, decimos entonces que $A$ está dotado de la estructura de un
\emph{$S$--$\cat{C}$-módulo filtrado} sobre el anillo filtrado $S$.

Al pasar a los objetos graduados asociados, todo $h_t$, para
$t\in\mathrm{F}^i(S)$, define un endomorfismo graduado
$\overline{h}_t$ de $\mathrm{gr}^\bullet(A)$, homogéneo de grado
$i$.
Además, este endomorfismo depende únicamente de la clase de $t$ en
$\mathrm{gr}^i(S)$, y se obtiene así un homomorfismo de
\emph{anillos graduados}
\[
  \overline{h}:
  \mathrm{gr}^\bullet(S)
  \longrightarrow
  \Hom_{\cat{C}}^\bullet
  \bigl(\mathrm{gr}^\bullet(A),\mathrm{gr}^\bullet(A)\bigr),
\]
donde el miembro derecho es el anillo de endomorfismos graduados de
$\mathrm{gr}^\bullet(A)$.
Decimos que $\mathrm{gr}^\bullet(A)$ está dotado de la estructura
de un \emph{$\mathrm{gr}^\bullet(S)$--$\cat{C}$-módulo graduado}.

Se sigue de \sref{0.13.6.5} que, para
$1\leq r\leq+\infty$, todo
$\overline{t}\in\mathrm{gr}^i(S)$ define canónicamente, sobre los
objetos bigraduados
\[
  \bigl(\mathrm{B}_r^{pq}(A)\bigr)_{(p,q)\in\bb{Z}\times\bb{Z}},
  \qquad
  \bigl(\mathrm{Z}_r^{pq}(A)\bigr)_{(p,q)\in\bb{Z}\times\bb{Z}},
\]
y
\[
  \mathrm{E}_r(A)
  =
  \bigl(\mathrm{E}_r^{pq}(A)\bigr)_{(p,q)\in\bb{Z}\times\bb{Z}},
\]
endomorfismos bigraduados de grado $(s,-s)$.
Sobre $\mathrm{E}_r(A)$, para $r$ finito, este endomorfismo conmuta con
el endomorfismo bigraduado definido por los $d_r^{pq}$.
Estos endomorfismos satisfacen las condiciones usuales de asociatividad y
distributividad respecto de la adición en
$\mathrm{gr}^\bullet(S)$ y en los objetos bigraduados en cuestión.
Para abreviar, decimos que estos últimos son
\emph{$\mathrm{gr}^\bullet(S)$--$\cat{C}'$-módulos bigraduados};
es inmediato que los $\alpha_r^{pq}$ definen un isomorfismo para
este tipo de estructura.

Para todo entero $n$, sea $\mathrm{B}_r^{(n)}(A)$
(respectivamente, $\mathrm{Z}_r^{(n)}(A)$ y
$\mathrm{E}_r^{(n)}(A)$) el subobjeto graduado de
$\mathrm{B}_r^{\bullet,\bullet}(A)$
(respectivamente, $\mathrm{Z}_r^{\bullet,\bullet}(A)$ y
$\mathrm{E}_r^{\bullet,\bullet}(A)$) constituido por los términos para
los que $p+q=n$ ($1\leq r\leq+\infty$).
Es inmediato que estos son
\emph{$\mathrm{gr}^\bullet(S)$--$\cat{C}'$-módulos graduados}.

Finalmente, todo $\overline{t}\in\mathrm{gr}^i(S)$ define, para todo
$n$, un endomorfismo graduado de grado $i$ sobre el objeto graduado
$\mathrm{gr}^\bullet(\RR^n\mathrm{T}(A))$, que adquiere así la
estructura de un
$\mathrm{gr}^\bullet(S)$--$\cat{C}'$-módulo graduado.
Para $p+q=n$, los $\beta^{pq}$ definen un isomorfismo
\[
  \mathrm{E}^{(n)}(A)
  \isoto
  \mathrm{gr}^\bullet(\RR^n\mathrm{T}(A))
\]
para este tipo de estructura.

Observemos que, cuando $\cat{C}'$ es la categoría de grupos abelianos, las
estructuras de $S$--$\cat{C}'$-módulo (respectivamente, de
$\mathrm{gr}^\bullet(S)$--$\cat{C}'$-módulo graduado o bigraduado) son
simplemente las estructuras usuales de $S$-módulo (respectivamente, de
$\mathrm{gr}^\bullet(S)$-módulo graduado o bigraduado).
\end{env}

\subsection{Funtores derivados de un límite proyectivo de argumentos}
\label{subsection:0.13.7}
\label{0.13.7}

\begin{env}[13.7.1]
\label{0.13.7.1}
Sean $\cat{C}$ y $\cat{C}'$ dos categorías abelianas, y supongamos
que todo objeto de $\cat{C}$ es isomorfo a un subobjeto de un
objeto inyectivo.
Sea
\[
  \mathrm{T}:\cat{C}\longrightarrow\cat{C}'
\]
un funtor covariante aditivo.
Consideremos un sistema proyectivo estricto
$\bb{A}=(A_k)_{k\in\bb{Z}}$ en $\cat{C}$, acotado inferiormente; más
precisamente, supongamos que $A_k=0$ para $k<k_0$.
A este sistema asociamos canónicamente, sobre cada $A_k$, la
filtración $(\mathrm{F}^p(A_k))_{p\in\bb{Z}}$ dada por
\sref{0.13.4.1.1}.
Como el sistema proyectivo es estricto, los morfismos canónicos
\begin{equation}
\label{0.13.7.1.1}
\tag{13.7.1.1}
  \mathrm{F}^i(A_h)/\mathrm{F}^j(A_h)
  \longrightarrow
  \mathrm{F}^i(A_k)/\mathrm{F}^j(A_k)
  \qquad(h\geq k)
\end{equation}
son isomorfismos para $i\leq j\leq k+1$.
Recordemos también que
$\mathrm{F}^{k_0}(A_k)=A_k$ y
$\mathrm{F}^{k+1}(A_k)=0$ para todo $k$.
\end{env}

\begin{env}[13.7.2]
\label{0.13.7.2}
Para cada $k$, construyamos una resolución inyectiva
$\mathrm{X}_k^\bullet=(\mathrm{X}_k^j)_{j\geq0}$ de $A_k$ que tenga las
propiedades de \sref{0.13.6.2}.
Los morfismos canónicos $A_{k+1}\longrightarrow A_k$ nos permiten, por
\sref{0.13.6.3}, definir para cada $k$ un morfismo de complejos
\[
  \mathrm{X}_{k+1}^\bullet\longrightarrow\mathrm{X}_k^\bullet
\]
compatible con las filtraciones y que hace de
$\bb{X}^\bullet=(\mathrm{X}_k^\bullet)_{k\in\bb{Z}}$ un sistema
proyectivo de complejos.
\oldpage[0\textsubscript{III}]{76}
Podemos suponer además que este sistema proyectivo es estricto.
En efecto, por \sref{0.13.7.1.1},
\[
  A_k
  \simeq
  A_{k+1}/\mathrm{F}^{k+1}(A_{k+1}).
\]
Así, al construir $\mathrm{X}_{k+1}^\bullet$, podemos tomar la
resolución inyectiva de
$A_{k+1}/\mathrm{F}^{k+1}(A_{k+1})$ \emph{igual} a
$\mathrm{X}_k^\bullet$.
Se sigue entonces de (M, V, 2.3) que el morfismo de complejos
$\mathrm{X}_{k+1}^\bullet\longrightarrow\mathrm{X}_k^\bullet$ puede
construirse de modo que, al pasar a los cocientes, induzca un
isomorfismo
\[
  \mathrm{X}_{k+1}^\bullet/
  \mathrm{F}^{k+1}(\mathrm{X}_{k+1}^\bullet)
  \simto
  \mathrm{X}_k^\bullet
\]
que respeta las filtraciones; esta es precisamente la condición de
\sref{0.13.5.1}.
\end{env}

\begin{env}[13.7.3]
\label{0.13.7.3}
Por construcción, el sistema proyectivo
$\mathrm{T}(\bb{X}^\bullet)$ de complejos
$\mathrm{T}(\mathrm{X}_k^\bullet)$ satisface las hipótesis de
\sref{0.13.5.3}.
Los resultados de \sref{0.13.5.4}, \sref{0.13.5.5} y
\sref{0.13.5.6} se aplican, por tanto, a las sucesiones espectrales
\[
  \mathrm{E}\bigl(\mathrm{T}(\mathrm{X}_k^\bullet)\bigr)
  =
  \mathrm{E}(A_k).
\]
Para $1\leq r\leq+\infty$, escribimos
$\mathrm{E}_r^{pq}(\bb{A})$ en lugar de
$\mathrm{E}_r^{pq}(\mathrm{T}(\bb{X}^\bullet))$
(cf.\ \sref{0.13.5.7} cuando $r=+\infty$), y utilizamos la misma convención
para la notación análoga.
En particular,
\begin{equation}
\label{0.13.7.3.1}
\tag{13.7.3.1}
  \mathrm{E}_1^{pq}(\bb{A})
  =
  \RR^{p+q}\mathrm{T}\bigl(\mathrm{gr}^p(\bb{A})\bigr),
\end{equation}
por \sref{0.13.6.4.2} y el hecho de que el sistema proyectivo
$(\mathrm{gr}^p(A_k))_{k\in\bb{Z}}$ es esencialmente constante.
Estos resultados, junto con \sref{0.13.4.3}, dan la proposición
siguiente, demostrada por primera vez por Shih Weishu mediante un método
diferente (inédito).
\end{env}

\begin{proposition}[13.7.4]
\label{0.13.7.4}
\emph{(Shih).}
Sea $n$ un entero.
Las dos condiciones siguientes son equivalentes:
\begin{enumerate}
  \item[{\rm a)}] Para todo par $(p,q)$ tal que $p+q=n$, el
  sistema proyectivo
  $(\mathrm{E}_r^{pq}(\bb{A}))_{r\geq2}$ es esencialmente constante.
  \item[{\rm b)}] El sistema proyectivo
  \[
    \RR^n\mathrm{T}(\bb{A})
    =
    \bigl(\RR^n\mathrm{T}(A_k)\bigr)_{k\in\bb{Z}}
  \]
  satisface {\rm(ML)}.
\end{enumerate}
Además, cuando se cumplen estas condiciones, existe un
isomorfismo canónico
\begin{equation}
\label{0.13.7.4.1}
\tag{13.7.4.1}
  \mathrm{gr}^p\bigl(\RR^n\mathrm{T}(\bb{A})\bigr)
  \isoto
  \mathrm{E}_\infty^{p,n-p}(\bb{A})
  \qquad(p\in\bb{Z}).
\end{equation}
\end{proposition}

\begin{proof}
Por \sref{0.13.5.6}, la condición {\rm a)} dice precisamente que, para
$p+q=n$, el sistema proyectivo
$(\mathrm{E}_\infty^{pq}(A_k))_{k\in\bb{Z}}$ es esencialmente
constante.
Por otra parte,
$\mathrm{gr}^p(\RR^n\mathrm{T}(A_k))$ es canónicamente isomorfo a
$\mathrm{E}_\infty^{p,n-p}(A_k)$.
Se sigue de \sref{0.13.4.3} que {\rm a)} y {\rm b)} son
equivalentes, y \sref{0.13.7.4.1} no es sino
\sref{0.13.5.6.1} en el caso considerado.
\end{proof}

\begin{corollary}[13.7.5]
\label{0.13.7.5}
Sea $\sh{F}=(\sh{F}_k)_{k\in\bb{N}}$ un sistema proyectivo de
haces de grupos abelianos que satisface las condiciones {\rm(i)}, {\rm(ii)}
y {\rm(iii)} de \sref{0.13.3.1}, y pongamos
$\sh{F}=\varprojlim_k\sh{F}_k$.
Supongamos que, para el funtor
$\sh{G}\longmapsto\Gamma(X,\sh{G})$, el sistema proyectivo
$(\mathrm{E}_r^{pq}(\sh{F}))_r$ es esencialmente constante siempre que
$p+q=n$ o $p+q=n+1$.
Sobre $\HH^{n+1}(X,\sh{F})$, consideremos la filtración
\[
  \mathrm{F}^p\bigl(\HH^{n+1}(X,\sh{F})\bigr)
  =
  \Ker\bigl(
    \HH^{n+1}(X,\sh{F})
    \longrightarrow
    \HH^{n+1}(X,\sh{F}_{p-1})
  \bigr).
\]
Entonces existe un isomorfismo canónico
\begin{equation}
\label{0.13.7.5.1}
\tag{13.7.5.1}
  \mathrm{gr}^p\bigl(\HH^{n+1}(X,\sh{F})\bigr)
  \isoto
  \mathrm{E}_\infty^{p,n-p+1}(\sh{F})
  \qquad(p\in\bb{Z}).
\end{equation}
\end{corollary}

\begin{proof}
\oldpage[0\textsubscript{III}]{77}
Apliquemos \sref{0.13.7.4} tomando como $\cat{C}$ la categoría de haces de
grupos abelianos sobre $X$, como $\cat{C}'$ la categoría de grupos abelianos,
y $\mathrm{T}=\Gamma$.
Da, para todo $p\in\bb{Z}$, un isomorfismo canónico
\[
  \mathrm{gr}^p\bigl(\RR^{n+1}\Gamma(\sh{F})\bigr)
  \isoto
  \mathrm{E}_\infty^{p,n-p+1}(\sh{F}).
\]
Además, por \sref{0.13.7.4}, el sistema proyectivo
$(\HH^n(X,\sh{F}_k))_{k\in\bb{Z}}$ satisface {\rm(ML)}; por tanto
\sref{0.13.3.1} da un isomorfismo canónico
\begin{equation}
\label{0.13.7.5.1b}
\tag{13.7.5.1}
  \HH^{n+1}(X,\sh{F})
  \isoto
  \varprojlim_k\HH^{n+1}(X,\sh{F}_k).
\end{equation}

Como el sistema proyectivo
$\RR^{n+1}\Gamma(\sh{F})$ satisface {\rm(ML)} por
\sref{0.13.7.4}, existe un isomorfismo canónico
\[
  \mathrm{gr}^p\bigl(\RR^{n+1}\Gamma(\sh{F})\bigr)
  \isoto
  \varprojlim_k
  \mathrm{gr}^p\bigl(\HH^{n+1}(X,\sh{F}_k)\bigr)
\]
por \sref{0.13.4.3}, y un isomorfismo canónico
\[
  \varprojlim_k
  \mathrm{gr}^p\bigl(\HH^{n+1}(X,\sh{F}_k)\bigr)
  \isoto
  \mathrm{gr}^p\left(
    \varprojlim_k\HH^{n+1}(X,\sh{F}_k)
  \right)
\]
por \sref{0.13.4.5}.
Solo queda comprobar que el isomorfismo precedente es
compatible con las filtraciones de sus dos lados.
Esto se sigue inmediatamente de las definiciones y, para todo $p$,
de la conmutatividad de
\[
  \xymatrix{
    \HH^{n+1}(X,\sh{F})
      \ar[r]^-{\sim}\ar[dr]
    &
    \varprojlim_k\HH^{n+1}(X,\sh{F}_k)
      \ar[d]
    \\
    &
    \HH^{n+1}(X,\sh{F}_{p-1}).
  }
\]
\end{proof}

\begin{env}[13.7.6]
\label{0.13.7.6}
Sea $S$ un anillo dotado de una filtración
$(\mathrm{F}^i(S))_{i\in\bb{Z}}$ tal que
$\mathrm{F}^0(S)=S$ (y por tanto
$\mathrm{gr}^i(S)=0$ para $i<0$).
Supongamos que cada $A_k$, con la filtración definida en
\sref{0.13.7.1}, es un
$S$--$\cat{C}$-módulo filtrado en el sentido de \sref{0.13.6.6}, y que los
morfismos $A_h\longrightarrow A_k$ para $k\leq h$ son morfismos de
$S$--$\cat{C}$-módulos filtrados.
Para abreviar, decimos entonces que $\bb{A}$ es un
\emph{sistema proyectivo de $S$--$\cat{C}$-módulos filtrados}.

Es inmediato que, para $k\leq h$ y
$1\leq r\leq+\infty$, los morfismos
\[
  \mathrm{B}_r^{pq}(A_h)\longrightarrow\mathrm{B}_r^{pq}(A_k),
  \qquad
  \mathrm{Z}_r^{pq}(A_h)\longrightarrow\mathrm{Z}_r^{pq}(A_k)
\]
son morfismos de
$\mathrm{gr}^\bullet(S)$--$\cat{C}'$-módulos bigraduados
(\sref{0.13.6.5}), y que, para $r$ finito, las familias
\[
  \bigl(\mathrm{Z}_r^{pq}(\bb{A})\bigr),\qquad
  \bigl(\mathrm{B}_r^{pq}(\bb{A})\bigr),\qquad
  \bigl(\mathrm{E}_r^{pq}(\bb{A})\bigr)
\]
son $\mathrm{gr}^\bullet(S)$--$\cat{C}'$-módulos
bigraduados, siendo los dos primeros submódulos de
$(\mathrm{E}_1^{pq}(\bb{A}))$.
Escribimos de nuevo
$\mathrm{Z}_r^{(n)}(\bb{A})$,
$\mathrm{B}_r^{(n)}(\bb{A})$ y
$\mathrm{E}_r^{(n)}(\bb{A})$ para los respectivos subobjetos graduados
obtenidos conservando únicamente los términos con $p+q=n$; estos son
$\mathrm{gr}^\bullet(S)$--$\cat{C}'$-módulos graduados.

Cuando el sistema proyectivo
$(\mathrm{E}_r^{pq}(\bb{A}))_r$ es esencialmente constante,
$\mathrm{E}_\infty^{pq}(\bb{A})$ es asimismo un
$\mathrm{gr}^\bullet(S)$--$\cat{C}'$-módulo bigraduado, y cada
$\mathrm{E}_\infty^{(n)}(\bb{A})$ es un
$\mathrm{gr}^\bullet(S)$--$\cat{C}'$-módulo graduado.
Además, para cada $k$, los isomorfismos
\[
  \beta^{p,n-p}:
  \mathrm{E}_\infty^{p,n-p}(A_k)
  \isoto
  \mathrm{gr}^p\bigl(\RR^n\mathrm{T}(A_k)\bigr)
\]
forman un isomorfismo de
$\mathrm{gr}^\bullet(S)$--$\cat{C}'$-módulos graduados de
$\mathrm{E}_\infty^{(n)}(A_k)$ sobre
$\mathrm{gr}^\bullet(\RR^n\mathrm{T}(A_k))$.
Bajo las hipótesis precedentes, el isomorfismo
\[
  \beta^{p,n-p}:
  \mathrm{E}_\infty^{(n)}(\bb{A})
  \isoto
  \varprojlim_k
  \mathrm{gr}^\bullet\bigl(\RR^n\mathrm{T}(A_k)\bigr)
\]
respeta, por tanto, también estas estructuras.
Lo mismo vale manifiestamente para el isomorfismo canónico
\[
  \mathrm{gr}^\bullet\bigl(\RR^n\mathrm{T}(\bb{A})\bigr)
  \isoto
  \varprojlim_k
  \mathrm{gr}^\bullet\bigl(\RR^n\mathrm{T}(A_k)\bigr);
\]
por consiguiente, los isomorfismos \sref{0.13.7.4.1} forman un
isomorfismo de
$\mathrm{gr}^\bullet(S)$--$\cat{C}'$-módulos graduados.
\end{env}

\begin{proposition}[13.7.7]
\label{0.13.7.7}
Sea $S$ un anillo noetheriano $\mathfrak{J}$-ádico.
Supongamos que $\cat{C}$ es una categoría abeliana en la que todo objeto
es isomorfo a un subobjeto de un objeto inyectivo, y sea
$\mathrm{T}$ un funtor covariante aditivo de $\cat{C}$ en la
categoría de grupos abelianos.
Sea $\bb{A}=(A_k)_{k\in\bb{Z}}$ un
\oldpage[0\textsubscript{III}]{78}
sistema proyectivo estricto de $S$--$\cat{C}$-módulos filtrados, para la
filtración $\mathfrak{J}$-ádica sobre $S$, acotado inferiormente.
Para un entero fijo $n$, supongamos que se cumple la condición siguiente:
\begin{itemize}
  \item[$(\mathrm{F}_n)$]
  \emph{El $\mathrm{gr}^\bullet(S)$-módulo graduado}
  \[
    \mathrm{E}_1^{(m)}(\bb{A})
    =
    \bigl(
      \RR^m\mathrm{T}(\mathrm{gr}^p(\bb{A}))
    \bigr)_{p\in\bb{Z}}
  \]
  \emph{de \sref{0.13.7.3.1} es de tipo finito para
  $m=n$ y $m=n+1$.}
\end{itemize}
Entonces:
\begin{enumerate}
  \item[{\rm(i)}] Los sistemas proyectivos
  $(\RR^n\mathrm{T}(A_k))_{k\in\bb{Z}}$ y
  $(\RR^{n+1}\mathrm{T}(A_k))_{k\in\bb{Z}}$ satisfacen {\rm(ML)}.
  \item[{\rm(ii)}] Si
  \[
    \RR'^{\,n}\mathrm{T}(\bb{A})
    =
    \varprojlim_k\RR^n\mathrm{T}(A_k),
  \]
  entonces $\RR'^{\,n}\mathrm{T}(\bb{A})$ es un $S$-módulo de tipo
  finito.
  \item[{\rm(iii)}] La filtración sobre
  $\RR'^{\,n}\mathrm{T}(\bb{A})$ definida por
  \[
    \mathrm{F}^p\bigl(\RR'^{\,n}\mathrm{T}(\bb{A})\bigr)
    =
    \Ker\bigl(
      \RR'^{\,n}\mathrm{T}(\bb{A})
      \longrightarrow
      \RR^n\mathrm{T}(A_{p-1})
    \bigr)
    \qquad(p\in\bb{Z})
  \]
  es $\mathfrak{J}$-buena; es decir,
  \[
    \mathfrak{J}\,
    \mathrm{F}^p\bigl(\RR'^{\,n}\mathrm{T}(\bb{A})\bigr)
    \subset
    \mathrm{F}^{p+1}\bigl(\RR'^{\,n}\mathrm{T}(\bb{A})\bigr)
  \]
  para todo $p$, con igualdad para todo $p$ suficientemente grande.
  En particular, la topología sobre
  $\RR'^{\,n}\mathrm{T}(\bb{A})$ definida por esta filtración es la
  topología $\mathfrak{J}$-ádica.
  \item[{\rm(iv)}] El sistema proyectivo
  $(\mathrm{E}_r^{pq}(\bb{A}))_r$ es esencialmente constante siempre que
  $p+q=n$ o $p+q=n+1$.
  Así, $\mathrm{E}_\infty^{pq}(\bb{A})$ está definido
  (\sref{0.13.5.7}), y los isomorfismos
  \begin{equation}
  \label{0.13.7.7.1}
  \tag{13.7.7.1}
    \mathrm{gr}^p\bigl(\RR'^{\,n}\mathrm{T}(\bb{A})\bigr)
    \isoto
    \mathrm{E}_\infty^{p,n-p}(\bb{A})
    \qquad(p\in\bb{Z})
  \end{equation}
  forman un isomorfismo de
  $\mathrm{gr}^\bullet(S)$-módulos graduados.
\end{enumerate}
\end{proposition}

\begin{proof}
En vista de \sref{0.13.7.7.1}, y teniendo en cuenta los
isomorfismos \sref{0.13.7.4.1}, por abuso de notación escribiremos también
$\RR^n\mathrm{T}(\bb{A})$ para el límite proyectivo
$\RR'^{\,n}\mathrm{T}(\bb{A})$ del sistema proyectivo
$(\RR^n\mathrm{T}(A_k))$.

El anillo graduado $\mathrm{gr}^\bullet(S)$ es noetheriano
(Bourbaki, \emph{Alg.\ comm.}, cap.~III, \textsection2,
n\textsuperscript{o}~9, cor.~5 del teorema~2).
Por tanto, la sucesión creciente de
$\mathrm{gr}^\bullet(S)$-submódulos graduados
$\mathrm{B}_r^{(m)}(\bb{A})$ de
$\mathrm{E}_1^{(m)}(\bb{A})$ es estacionaria para $m=n$ y
$m=n+1$.
La condición {\rm(b)} de \sref{0.11.1.10} queda, por tanto, satisfecha.
Se sigue que la condición {\rm(a)} de \sref{0.13.7.4} se cumple para
$n$ y para $n+1$, lo que demuestra {\rm(i)}.

Además, los isomorfismos \sref{0.13.7.4.1}, junto con las
observaciones de \sref{0.13.7.6}, muestran que
$\mathrm{gr}^\bullet(\RR^n\mathrm{T}(\bb{A}))$ es un
$\mathrm{gr}^\bullet(S)$-módulo graduado isomorfo a
\[
  \mathrm{E}_\infty^{(n)}(\bb{A})
  =
  \mathrm{Z}_\infty^{(n)}(\bb{A})
  /
  \mathrm{B}_\infty^{(n)}(\bb{A}).
\]
Como $\mathrm{Z}_\infty^{(n)}(\bb{A})$ es un submódulo de
$\mathrm{E}_1^{(n)}(\bb{A})$, es de tipo finito, y también lo es
$\mathrm{E}_\infty^{(n)}(\bb{A})$.
Finalmente, para la filtración
$(\mathrm{F}^p(\RR'^{\,n}\mathrm{T}(\bb{A})))$, se sigue de
\sref{0.13.4.5} que
\[
  \mathrm{gr}^\bullet\bigl(\RR^n\mathrm{T}(\bb{A})\bigr)
  \quad\hbox{y}\quad
  \mathrm{gr}^\bullet\bigl(\RR'^{\,n}\mathrm{T}(\bb{A})\bigr)
\]
son $\mathrm{gr}^\bullet(S)$-módulos graduados isomorfos.
Esto demuestra {\rm(iv)}.
Las aserciones {\rm(ii)} y {\rm(iii)} se seguirán de estos resultados
y del lema siguiente.
\end{proof}

\begin{lemma}[13.7.7.2]
\label{0.13.7.7.2}
Sea $S$ un anillo noetheriano $\mathfrak{J}$-ádico, y sea $M$
un $S$-módulo dotado de una filtración codiscreta
$(\mathrm{F}^p(M))_{p\in\bb{Z}}$ tal que
\[
  \mathfrak{J}\mathrm{F}^p(M)
  \subset
  \mathrm{F}^{p+1}(M).
\]
Así, $M$ es un módulo filtrado sobre $S$ dotado de su
filtración $\mathfrak{J}$-ádica.
Supongamos además que $M$ es separado para la topología definida por
$(\mathrm{F}^p(M))$.
Entonces las condiciones siguientes son equivalentes:
\begin{enumerate}
  \item[{\rm a)}] $M$ es un $S$-módulo de tipo finito y
  $(\mathrm{F}^p(M))$ es una filtración $\mathfrak{J}$-buena.
  \item[{\rm b)}] $\mathrm{gr}^\bullet(M)$ es un
  $\mathrm{gr}^\bullet(S)$-módulo de tipo finito.
  \item[{\rm c)}] Los $\mathrm{gr}^p(M)$ son $S$-módulos de
  tipo finito y, para todo $p$ suficientemente grande, los homomorfismos
  canónicos
  \begin{equation}
  \label{0.13.7.7.3}
  \tag{13.7.7.3}
    \mathfrak{J}\otimes_S\mathrm{gr}^p(M)
    \longrightarrow
    \mathrm{gr}^{p+1}(M)
  \end{equation}
  \oldpage[0\textsubscript{III}]{79}
  son sobreyectivos.
  Están inducidos por
  $\mathfrak{J}\otimes_S\mathrm{F}^p(M)
  \longrightarrow\mathrm{F}^{p+1}(M)$, teniendo en cuenta el hecho
  de que la imagen del homomorfismo compuesto
  \[
    \mathfrak{J}\otimes_S\mathrm{F}^{p+1}(M)
    \longrightarrow
    \mathfrak{J}\otimes_S\mathrm{F}^p(M)
    \longrightarrow
    \mathrm{F}^{p+1}(M)
  \]
  es
  $\mathfrak{J}\mathrm{F}^{p+1}(M)
  \subset\mathrm{F}^{p+2}(M)$.
\end{enumerate}

Para la demostración, véase Bourbaki, \emph{Alg.\ comm.}, cap.~III,
\textsection3, n\textsuperscript{o}~1, prop.~3.
\end{lemma}

\begin{env}[13.7.7.4]
\label{0.13.7.7.4}
Para aplicar el lema \sref{0.13.7.7.2}, queda observar que la
topología sobre $\RR'^{\,n}\mathrm{T}(\bb{A})$ definida por la filtración
considerada lo convierte en un $S$-módulo separado y completo:
esta es la topología de límite proyectivo de los grupos discretos
$\RR^n\mathrm{T}(A_k)$.
Por otra parte, si $A_k=0$ para $k<k_0$, entonces también
$\RR^n\mathrm{T}(A_k)=0$ para $k<k_0$; por tanto
\[
  \mathrm{F}^{k_0}\bigl(\RR'^{\,n}\mathrm{T}(\bb{A})\bigr)
  =
  \RR'^{\,n}\mathrm{T}(\bb{A}).
\]
Así, todas las hipótesis del lema quedan efectivamente satisfechas.
\end{env}

\begin{corollary}[13.7.8]
\label{0.13.7.8}
Si se cumple la hipótesis $(\mathrm{F}_n)$, entonces para todo $f\in S$
existe un isomorfismo canónico
\begin{equation}
\label{0.13.7.8.1}
\tag{13.7.8.1}
  \varprojlim_k\bigl(\RR^n\mathrm{T}(A_k)\bigr)_f
  \isoto
  \RR^n\mathrm{T}(\bb{A})\otimes_S S_{\{f\}}.
\end{equation}
\end{corollary}

\begin{proof}
En efecto, $\RR^n\mathrm{T}(\bb{A})$ es un $S$-módulo de tipo finito,
y $S_{\{f\}}$ es una $S$-álgebra ádica noetheriana
(\sref[0\textsubscript{I}]{0.7.6.11}), a saber, la completación separada
de $S_f$ para la topología $\mathfrak{J}$-preádica
(\sref[0\textsubscript{I}]{0.7.6.2}).
Se sigue de \sref[0\textsubscript{I}]{0.7.7.8} y
\sref[0\textsubscript{I}]{0.7.7.1} que
\[
  \RR^n\mathrm{T}(\bb{A})\otimes_S S_{\{f\}}
\]
es isomorfo a la completación separada de
$\RR^n\mathrm{T}(\bb{A})\otimes_S S_f$ para la
topología $\mathfrak{J}$-preádica.
Un sistema fundamental de entornos de cero para esta topología
es
\[
  \bigl(
    \mathfrak{J}^p\RR^n\mathrm{T}(\bb{A})
  \bigr)\otimes_S S_f;
\]
los grupos
\[
  \mathrm{F}^p\bigl(\RR^n\mathrm{T}(\bb{A})\bigr)
  \otimes_S S_f
\]
forman, por tanto, también un sistema tal.
Este último grupo es el núcleo de la aplicación canónica
\[
  \bigl(\RR^n\mathrm{T}(\bb{A})\bigr)_f
  \longrightarrow
  \bigl(\RR^n\mathrm{T}(A_{p-1})\bigr)_f.
\]
Por consiguiente, el grupo separado asociado a
$\RR^n\mathrm{T}(\bb{A})\otimes_S S_f$ se identifica con un subgrupo
$G$ de
\[
  \varprojlim_k\bigl(\RR^n\mathrm{T}(A_k)\bigr)_f.
\]
Pero el sistema proyectivo
$((\RR^n\mathrm{T}(A_k))_f)_k$ satisface manifiestamente {\rm(ML)}, y la
imagen de $(\RR^n\mathrm{T}(\bb{A}))_f$ en cada
$(\RR^n\mathrm{T}(A_k))_f$ es igual a la imagen común de los
$(\RR^n\mathrm{T}(A_h))_f$ para todo $h\geq k$ suficientemente grande.
Se sigue inmediatamente que $G$ es \emph{denso en todas partes} en
$\varprojlim_k(\RR^n\mathrm{T}(A_k))_f$.
Como este último grupo es separado y completo, queda demostrado el
corolario.
\end{proof}
